\documentclass[11pt]{article}
\usepackage[margin=1in]{geometry}
\usepackage{amsmath,amssymb,booktabs,longtable}
\usepackage[hidelinks]{hyperref}
\newcommand{\dd}{\mathrm d}
\newcommand{\eq}{\mathrm{eq}}
\title{Pair conversion changes the cold-carrier abundance\\
A reciprocal expanding-gas calculation and its limits}
\author{Zimmerman programme: CPB-PC-2026-10-16}
\date{16 October 2026}
\begin{document}
\maketitle
\begin{abstract}
The scalar carrier contains a pole-free pair-conversion reaction that must
be included alongside source decay and inverse decay. Their combination
changes the eventual daughter count while conserving charge and energy.
We derive its tree cross section and thermal rate, and execute an expanding
Maxwell--Boltzmann gas comparison using fast source partial equilibrium.
One specified initial state nearly doubles its weighted daughter count;
a hotter comparison fails the preset dilution gate and is retained as such.
The evolution supplies concrete kinetic evidence, not a primordial production
history, a relic-density fit or a closed gravitational theory. A pressure-aware
generalization of the cold--vacuum exchange identity separates internal
rest-mass conversion from exchange with dark energy.
\end{abstract}

\section{Exact target and provenance}
The full user goal remains a physical completion of the theory, a defensible
particle bridge, and a complete account of cold and dark energy. This
checkpoint advances the missing collision/abundance calculation. It does
not replace that goal with a chosen initial-value problem.

The previous goal turn was progress: it established a thermal obstruction,
a sufficient alternative scalar branch, and a mixed-charge six-point
threshold amplitude. This turn reads the actual previous records before
extending the collision system. The checkout remains
\texttt{370dfb58006e330f210f5e306297e822188e37ba}, with shared dirty state.
Inputs and software are pinned by
\texttt{pair\_conversion\_run\_004/manifest.json}. The frozen theory and
other lanes remain unchanged. This is self-review.

We choose the separately specified $m=2$ MeV family member, not the earlier
thermally unstable 1-MeV point. It satisfies the previous sufficient
condition $M_0^2\ge g_c^2/\lambda_{\rm tot}$ within that thermal comparison.
The choices $u_*=.02$, $r_\lambda=2$ and
$\sigma_0/m=.3\ \mathrm{cm^2/g}$ remain inputs. The source mass is
$M=2m\sqrt{1+u_*^2}$, and $g_c,\lambda_{\rm tot}$ are the same derived
parameter values in the pinned mass-family record. Natural units apply.

\section{A missing reaction and an exact cross section}
The charge-preserving scalar interaction
\begin{equation}
 \mathcal L\supset-g_c(X^*\Phi^2+X\Phi^{*2})
 -\lambda_{\rm tot}|\Phi|^4
\end{equation}
allows
\begin{equation}
 \Phi+\Phi^*\ \rightleftarrows\ X+X^*.
\label{eq:pair}
\end{equation}
The cubic vertex is $-2ig_c$. For the forward process, one internal daughter
joins the two vertices; fixed charges give a single $t$-channel tree:
\begin{equation}
 \mathcal M=-\frac{4g_c^2}{t-m^2}.
\end{equation}
There is no crossed identical-final graph because $X$ and $X^*$ are distinct.
This amplitude is not the same as opposite-charge daughter elastic scattering.

At invariant $s\ge4M^2$ put
$p=\sqrt{s/4-m^2}$, $k=\sqrt{s/4-M^2}$ and
\begin{equation}
 t=m^2+M^2-s/2+2pk\cos\theta.
\end{equation}
Both initial and final energies are $\sqrt s/2$, so
$t=-|\mathbf p-\mathbf k|^2\le0$. The internal daughter is never on shell.
With $A=s/2-M^2$, $B=2pk$,
\begin{equation}
 A^2-B^2=M^4+m^2(s-4M^2)>0,\qquad
 \int_{-1}^1\frac{\dd z}{(A-Bz)^2}=\frac2{A^2-B^2}.
\end{equation}
Direct angular integration therefore gives
\begin{align}
 \sigma_f(s)&=\frac{g_c^4}{\pi s}
 \sqrt{\frac{s-4M^2}{s-4m^2}}
 \frac1{M^4+m^2(s-4M^2)},\\
 \sigma_r(s)&=\frac{g_c^4}{\pi s}
 \sqrt{\frac{s-4m^2}{s-4M^2}}
 \frac1{M^4+m^2(s-4M^2)}.
\end{align}
They obey
$(s-4m^2)\sigma_f=(s-4M^2)\sigma_r$.
There are no initial/final identical-pair factorials for either direction.
The reverse cross section diverges kinematically as $1/k$ at its threshold,
but its physical velocity product has the finite limit
\begin{equation}
 (\sigma_r v_{\rm M\o l})_0
 =\frac{g_c^4}{2\pi M^6}\sqrt{1-m^2/M^2}
 =9.50486\,10^{-8}\ \mathrm{GeV}^{-2}
\end{equation}
for the selected point. Angular quadrature and detailed-balance controls
check the normalization independently of the closed angular formula.

\section{Why a two-to-two reaction changes the final rest count}
Use a charge-symmetric state, with per-sign densities $n=n_\Phi=n_{\Phi^*}$
and $b=n_X=n_{X^*}$. Let $D$ be net source decay rate per sign and $P$
net forward pair rate. After visible production has ceased,
\begin{equation}
 \dot n+3Hn=2D-P,\qquad
 \dot b+3Hb=-D+P.
\end{equation}
The useful weighted count per sign is $w=n+2b$, the daughter count after
all remaining sources decay. It satisfies
\begin{equation}
 \boxed{\dot w+3Hw=P.}
\end{equation}
Forward pair conversion consumes one daughter of each sign and creates
one source of each sign; subsequent decays produce four daughters instead
of the initial two. The inverse cycle reduces the eventual daughter count.
The net conserved charge stays zero. Source decay alone conserves $w$;
the pair channel does not.

For a common Maxwell--Boltzmann kinetic temperature $T_d$, write fugacities
$z_\Phi,z_X$ relative to zero chemical potential. Detailed balance gives
\begin{equation}
 D=\gamma_d^\eq(z_X-z_\Phi^2),\qquad
 P=\gamma_p^\eq(z_\Phi^2-z_X^2),\qquad
 \gamma_p^\eq=(n_X^\eq)^2\langle\sigma_r v_{\rm M\o l}\rangle.
\end{equation}
Fast source decay/inverse decay imposes $z_X=z_\Phi^2$. Full chemical
equilibrium would additionally require $z_\Phi=1$; it is not imposed.
The collision structure, not dimensional analysis or the conserved charge,
determines whether and when that second equilibrium is approached.

The standard invariant reaction measure and flux definitions are checked
against De Romeri et al. equations (21)--(25) \cite{rates}. Our use of
classical distributions is an explicit reduction, not their Bose treatment.
For equal-mass distinguishable incoming sources the MB thermal average is
\begin{equation}
 \langle\sigma_rv_{\rm M\o l}\rangle=
 \frac{1}{8M^4T_d K_2(M/T_d)^2}
 \int_{4M^2}^\infty\sigma_r(s)(s-4M^2)\sqrt{s}
 K_1(\sqrt{s}/T_d)\dd s.
\end{equation}
We integrate with $y=(\sqrt s-2M)/T_d$ and exponentially scaled Bessels.
A generalized Laguerre weight $\sqrt y e^{-y}$ removes the threshold
square root. Across 25 temperatures in $T_d/m\in[10^{-5},10]$, 64- and
128-node rates differ by at most $3.45\,10^{-7}$ relative. The nonrelativistic
limit agrees with the displayed velocity product.

\section{Reciprocal energy closure without a rest-energy cancellation}
The gas is isolated from the visible sector during this comparison. Its
total dark energy and pressure satisfy
\begin{equation}
 \dot\epsilon_d+3H(\epsilon_d+p_d)=0.
\end{equation}
Pair production can increase rest mass only by reducing kinetic energy or
the other stored dark energy; it supplies no new total energy.
The thermodynamics uses an ideal quasiparticle equation of state. Interaction
and virial contributions to that equation of state are not included here.

For numerical stability define $\delta=M-2m$ and the reduced fugacity
$y=z_\Phi e^{-m/T_d}$. Write
\begin{align}
 A&=\frac{m^2T_d}{2\pi^2}e^{m/T_d}K_2(m/T_d),\\
 B&=\frac{M^2T_d}{2\pi^2}e^{M/T_d}K_2(M/T_d)e^{-\delta/T_d},\\
 n&=Ay,\quad b=By^2,\quad
 y=\frac{2w}{A+\sqrt{A^2+8Bw}}.
\end{align}
The last expression is the positive quadratic solution and stays well
conditioned for exponentially small equilibrium densities. With partial
source equilibrium the pair rate is
\begin{equation}
 P=B^2\langle\sigma_rv\rangle y^2
 [e^{-2m/T_d}-y^2].
\end{equation}
It can have either sign along the expanding trajectory.

For $e_i=m_iK_1(m_i/T_d)/K_2(m_i/T_d)+3T_d$, put
\begin{equation}
 k_\Phi=e_\Phi-m,\quad k_X=e_X-2m,\quad
 k=nk_\Phi+bk_X,\quad p=(n+b)T_d.
\end{equation}
Here $\epsilon_d=2(mw+k)$ and $p_d=2p$. Subtracting weighted rest energy
gives the exact moment identity within this closure,
\begin{equation}
 \dot k+3H(k+p)=-mP.
\end{equation}
If $P>0$, increasing the eventual rest count costs the kinetic/internal
budget. $P<0$ releases that rest energy into the remaining gas.

Let $c_i=\dd e_i/\dd T_d>0$. At fixed temperature and fixed weighted count,
respectively,
\begin{align}
 k_w&=\frac{nk_\Phi+2bk_X}{n+4b},\\
 k_T&=nc_\Phi+bc_X+
 \frac{nb(2k_\Phi-k_X)^2}{(n+4b)T_d^2}>0.
\end{align}
With $\ell=\ln a$ the actual two-variable evolution is
\begin{align}
 \frac{\dd\ln Y_w}{\dd\ell}&=\frac{P}{Hw},\\
 \frac{\dd T_d}{\dd\ell}&=
 \frac{-3p+3nb(k_X-2k_\Phi)/(n+4b)-(P/H)(m+k_w)}{k_T},\qquad
 Y_w=\frac{2w}{s_{\rm SM}}.
\end{align}
No rest-mass term is discarded. In the daughter-only elastic limit these
equations reduce to fixed $Y_w$ and $T_d\propto a^{-2}$ when nonrelativistic.
Thus rapid elastic scattering alone need not heat a homogeneous cold gas.
It changes the distribution and transport; pressure evolution follows its
energy moment. Number conversion changes that conclusion and is calculated
here explicitly.

There is also an independent thermodynamic check. The MB entropy density
per sign is
\begin{equation}
 s_d=\frac{k}{T_d}+n+b-w\ln y,\qquad
 \mu_\Phi=m+T_d\ln y,\quad \mu_X=2\mu_\Phi.
\end{equation}
The first law $\dd(mw+k)=T_d\dd s_d+\mu_\Phi\dd w$ and its Euler
identity give
\begin{equation}
 \dot s_d+3Hs_d=-\frac{\mu_\Phi}{T_d}P\ge0.
\end{equation}
Indeed $P=\gamma_p^\eq(z_\Phi^2-z_\Phi^4)$ and
$\mu_\Phi/T_d=\ln z_\Phi$ have opposite signs. Pair-off trajectories
conserve comoving entropy; pair-on trajectories produce it. This is a
conditional MB reaction-network entropy statement, not a claim about
uncomputed quantum kinetics.

\section{The executed expanding-gas comparison}
The initial visible temperature is 10 MeV and the final visible temperature
is 1 keV. The background contains a photon/electron ideal gas and three
massless neutrinos, with instantaneous neutrino decoupling at 2 MeV.
Electron Fermi integrals are included; visible entropy and the neutrino
temperature are evolved consistently. This is an inherited radiation
background, not the framework's native cosmology or precision neutrino
decoupling calculation. Its 500-row entropy table conserves the declared
entropy to the checked tolerance.

The specified initial weighted yield is
\begin{equation}
 Y_{w,i}=\frac{.12}{2.742\,10^8(m/\mathrm{GeV})}
 =2.18818\,10^{-7}.
\end{equation}
It is the old abundance-calibrated input, not a completed prediction from
the Higgs era. Initial source densities follow the assumed fast partial
equilibrium. Three independent initial kinetic temperatures are investigated;
none is claimed to have been produced by the earlier portal calculation.

\begin{center}
\begin{tabular}{r r r l}
\toprule Initial $T_d/m$ & Final $Y_w/Y_{w,i}$ & Final $T_d$ (GeV) & Dilution gate\\\midrule
.3, pair off & 1 & $5.59317\,10^{-12}$ & met\\
.3, pair on & 1.03195 & $9.58651\,10^{-12}$ & met\\
1, pair off & 1 & $4.21977\,10^{-11}$ & met\\
1, pair on & 1.97924 & $1.67032\,10^{-11}$ & met\\
3, pair off & 1 & $3.46622\,10^{-10}$ & met\\
3, pair on & 4.52622 & $3.39091\,10^{-11}$ & failed\\\bottomrule
\end{tabular}
\end{center}
Every number is a result for this reduced initial-value problem. In
particular, $1.97924$ is not an observed relic excess or a physical Higgs
production prediction. Turning the pair reaction off keeps the weighted
count exactly fixed but retains the same expansion/source closure.

The initial-energy ceilings on the yield ratios are 1.57214, 3.36915 and
9.15433, respectively. They follow from nonnegative pressure and
$\epsilon_d\ge2mw$: comoving energy cannot increase in the isolated gas.
The computed yield peaks 1.09428, 2.30842 and 6.02160 remain below those
ceilings. The subsequent decline is reciprocal reverse conversion and
freeze-out, rather than arbitrary deletion of mass.

For pair-on runs the independent expansion-work residual
\begin{equation}
 \frac{a_f^3\epsilon_{d,f}-\epsilon_{d,i}
 +\int_0^{\ell_f}3a^3p_d\dd\ell}{\epsilon_{d,i}}
\end{equation}
is below $5\,10^{-10}$ in magnitude. The weighted-number collision moment
is independently integrated as well. The refined $.3$ run has energy and
number residuals below $7\,10^{-12}$. ODE/rate refinement changes its final
yield by less than $10^{-9}$ relative. These numerical invariants verify
the implemented reduced equations, not the completeness of the collision set.
The pair-on final comoving entropy ratios are 1.12636, 1.62315 and 2.98564.
Their separately integrated entropy-production balances agree to better
than $1.4\,10^{-8}$ relative to initial entropy. Pair-off entropy ratios
equal one within $7\,10^{-10}$.

\section{Validity diagnostics and preserved failures}
The original script, pinned in failed run 001, required maximum MB
occupation $y<.02$, source-decay/evolution rate ratio $>10^5$, and a lower
bound on elastic/evolution rate ratio $>10$. These thresholds are retained.
The pair-on maximum occupations are .0059944, .0114825 and .0229111.
The hot case fails the original dilution gate. It remains an approximate
classical comparison and is not promoted to a result meeting that gate.
The failure means the declared small-occupation criterion is missed; it is
not evidence for a condensate or a universal failure of a classical expansion.

All three pair-on runs have elastic lower bounds at least 500 times their
largest local evolution scale under the chosen normalization, and source
decay/evolution ratios above $7.5\,10^8$. The elastic bound uses the
pole-free daughter opposite-charge inequality
\begin{equation}
 \sigma_{+-}(s)\ge\frac{(\lambda_{\rm tot}-g_c^2/M^2)^2}{\pi s}.
\end{equation}
These are quantitative support for the kinetic and fast-decay reductions,
not a solution of the full momentum-space collision operator. The
occupation/evolution checks are sampled on the stored 401-point trajectories;
no global-in-time certificate is claimed from those samples.

Runs 001 and 002 stopped on the hot-case dilution assertion; run 002 records
the measured diagnostics before stopping. Run 003 preserves the same
threshold, exports all computed trajectories and reports the failed scope
gate separately from numerical execution checks. It does not loosen the
gate to manufacture a pass. The final authoritative run 004 records
585 numerical/algebraic controls, including 500 background-entropy checks.
Run 003 remains the pre-entropy checkpoint. This count is not a
measure of physical completeness.

\section{What the full reaction system still needs}
The explicit pair process and source inverse decays already account for
real intermediate sources in a reaction network. At nonzero daughter
momenta, parts of the previous six-point amplitude can develop on-shell
source propagators and overlap those sequential processes. Adding a raw
finite-width $4\leftrightarrow2$ rate on top would count those histories twice.
The remaining collision hierarchy needs consistent real-intermediate-state
subtraction or an off-shell kinetic derivation. The general issue is
documented by Ala-Mattinen et al. \cite{ris}; their subtraction prescription
has not been transferred to this multiparticle scalar system here.

The preceding exact zero-incoming-momentum six-point result remains valid:
its intermediate source invariants lie below the positive-gap source pole.
It neither supplies the finite-temperature subtracted remainder nor bounds
it at the warmer initial states evolved in this report. Those omitted
rates have not been bounded sufficiently to establish the displayed yield
ratios as predictions of the complete scalar action.

Other remaining obligations are the real Higgs-era source and thermal
history, the resulting initial distributions, Bose factors and medium
dispersion/width corrections, and the full kinetic response where timescale
reductions fail. At sufficiently late nonrelativistic temperatures there
is no reason to add a fictitious visible-temperature dark thermal mass;
it must be calculated from the actual states.

\section{Cold energy exchange is not automatically dark-energy exchange}
Internal conversion here moves energy between particle rest mass and
kinetic energy while the total dark gas obeys its conservation equation.
It does not require the vacuum to provide or absorb energy. In a bound
galaxy, cosmological expansion is not a substitute for a physical local
formation receiver; the prior halo-settling obligation remains open.

The previous cold--vacuum consistency identity assumed a dust component.
For a component with nonzero $w_c=p_c/(\rho_cc^2)$, mass-equivalent total
energy density $\rho_c$, conserved baryons, and $Q$ positive from dark
energy into that component, the correct equations are
\begin{align}
 \dot\rho_c+3H(1+w_c)\rho_c&=Q,\\
 \dot\rho_{\rm DE}+3H(1+w_{\rm DE})\rho_{\rm DE}&=-Q,\\
 a_0^2&=\kappa^2c^2G\rho_{\rm DE}.
\end{align}
With constant $G,c,\kappa$, elimination of $Q$ gives
\begin{equation}
 \boxed{2\frac{\dd\ln a_0}{\dd\ln a}+3(1+w_{\rm DE})
 +\frac{\rho_c}{\rho_{\rm DE}}
 \left[\frac{\dd\ln(\rho_c/\rho_b)}{\dd\ln a}+3w_c\right]=0.}
\end{equation}
The proof is direct substitution of the two continuity equations after
differentiating the scale postulate. For this isolated gas, $Q=0$ even
when its eventual rest count changes. Treating $m$ times that count as
the entire $\rho_c$ during the warm conversion epoch would falsely infer
vacuum exchange. The old dust identity is recovered at $w_c=0$.
This distinction is necessary for a coherent cold/dark-energy story.
It derives neither the scale postulate nor $\kappa$, and fixes no vacuum
energy magnitude.

\section{Next route and completion status}
The executed result changes the next research decision: source production
cannot be converted into final daughter abundance by multiplying by four
and stopping. The pair channel, inverse decay and finite-temperature
irreducible number change must be evolved with reciprocal energy exchange.
The initial energy distribution can affect the final count substantially.

Next derive the real-intermediate-state-consistent irreducible scalar
collision kernel and its error relative to the explicitly evolved network.
Then replace the specified initial states with the stabilized candidate's
actual Higgs-era production and thermal history, and compute perturbations
from that same state. A portal recalibration before those steps would hide
the unresolved abundance mechanism rather than close it.

The galaxy-response action, finite counted halo formation, source partition,
lensing metric, $\kappa$, independent parameter evidence and common-model
empirical tests remain unresolved. The unique/novel-path ambition remains
subject to literature comparison: reaction-network reprocessing and dark
temperature evolution are known mechanisms. This report establishes their
specific consequences for the present chosen scalar comparison, not global
novelty or full closure.

\begin{thebibliography}{9}
\bibitem{rates} V. De Romeri, D. Karamitros, O. Lebedev and T. Toma,
\emph{Neutrino dark matter and the Higgs portal: improved freeze-in analysis},
arXiv:2003.12606v2, equations (21)--(25) and the discussion of thermal masses.
\url{https://arxiv.org/html/2003.12606v2}.
Only invariant rate/flux conventions and methodological overlap are used;
the quoted model's sterile-neutrino results are not imported.
\bibitem{ris} K. Ala-Mattinen, M. Heikinheimo, K. Kainulainen and K. Tuominen,
\emph{Anatomy of real intermediate state-subtraction scheme},
arXiv:2309.16615v2, sections I--III, especially the on-shell overlap in
the introduction and equations (1)--(4).
\url{https://arxiv.org/html/2309.16615v2}.
The source authenticates the double-counting issue; it does not supply a
completed multiparticle subtraction for this model.
\end{thebibliography}
\end{document}
